\chapter{مقدمه}
\label{ch:intro}

\section{بیان مسئله}
خلیج فارس، تنگه هرمز و دریای عمان گذرگاه بخش بزرگی از نفت و کالای دریایی جهان‌اند. برای
پایش این آبراه‌ها، داده اصلی امروز سامانه شناسایی خودکار (\lr{AIS}) است: هر شناورِ مشمول
مقررات، شناسه، موقعیت، سرعت و جهت خود را به‌طور منظم پخش می‌کند و گیرنده‌های ساحلی و
ماهواره‌ای این پیام‌ها را جمع می‌کنند. اما \lr{AIS} خوداظهاری است: شناور می‌تواند فرستنده را
خاموش کند، موقعیت ساختگی بفرستد یا در منطقه‌ای باشد که اخلال در ناوبری ماهواره‌ای
(\lr{GNSS}) موقعیتش را جابه‌جا می‌کند. شناوری که در دریا هست ولی در \lr{AIS} دیده نمی‌شود را
«شناور تاریک» می‌نامند.

تصاویر رادار روزنه مصنوعی (\lr{SAR}) راه مستقلی برای دیدن شناورهاست. رادار خود انرژی
می‌فرستد و بازتاب آن را می‌گیرد، پس در شب و از پشت ابر هم تصویر می‌گیرد، و بدنه فولادی
شناور در آن به‌صورت نقطه‌ای روشن روی دریای تیره دیده می‌شود، چه شناور \lr{AIS} داشته باشد چه
نه. ماهواره‌های \lr{Sentinel-1} برنامه کوپرنیک اروپا این تصاویر را به‌رایگان و با تفکیک حدود
۱۰ متر در دسترس می‌گذارند. پرسش این پژوهش این است که یک مدل یادگیری عمیق تا چه اندازه
می‌تواند شناورها را در این تصاویر بیابد، چگونه می‌توان آن را بدون برچسب‌گذاری دستی آموزش
داد، دقتش را در برابر مرجعی مستقل سنجید، و از کنار هم گذاشتن آشکارسازی‌های راداری با
\lr{AIS} نامزدهای شناور تاریک را با شواهدشان یافت.

\section{هدف‌ها}
\begin{enumerate}
\item آموزش یک آشکارساز یادگیری عمیق برای یافتن شناور در تصاویر \lr{Sentinel-1} بدون
برچسب‌گذاری دستی، که روی پردازنده یک رایانه همراه اجرا شود.
\item سنجش آن در برابر مرجعی مستقل، یعنی آشکارسازی‌های راداری منتشرشده
\lr{Global Fishing Watch} (\lr{GFW})، و مقایسه با روش کلاسیک \lr{CFAR} روی گذرهایی که در
آموزش و انتخاب مدل به کار نرفته‌اند.
\item آزمودن تعمیم مدل به زمانی دیگر، به ماهواره‌های تازه \lr{Sentinel-1C/1D} و به منطقه‌ای
دیگر (ساحل ایران).
\item یافتن نامزدهای شناور تاریک با تلفیق آشکارسازی‌ها با حضور ساعتی \lr{AIS} و ثبت شواهد
هر نامزد.
\item ساخت یک کنسول وب که هر آشکارسازی را همراه تصویر راداری‌اش و با بیان صادقانه سن داده و
دقت سنجیده‌شده نشان دهد.
\end{enumerate}

\section{دستاوردها}
\begin{itemize}
\item \textbf{برچسب‌گذاری خودکار:} برچسب‌های آموزشی از تطبیق آشکارسازی‌های \lr{CFAR} با
خانه‌های آشکارسازی \lr{GFW} ساخته شد (\dataChips{} برش تصویری از \dataScenes{} تصویر).
\item \textbf{سوگیری برچسب و رفع آن:} مدل نخست، که فقط شناورهای دارای \lr{AIS} را مثبت
می‌دید، شناورهای بدون \lr{AIS} را کمتر از \lr{CFAR} می‌یافت (بازیابی \darkRecallRunOne{} در
برابر \darkRecallCfar). افزودن \dataTrainDark{} آشکارسازی بدون \lr{AIS} به نمونه‌های مثبت
بازیابی را به \darkRecallModel{} رساند.
\item \textbf{سنجش نگه‌داشته:} روی ۹۲ گذر آزمون، \lr{F1} مدل نهایی \testFOneModel{} و
\lr{CFAR} \testFOneCfar{} بود؛ پیکربندی مدل فقط روی گذرهای اعتبارسنجی برگزیده شد.
\item \textbf{تعمیم:} مدل بدون بازتنظیم روی \lr{Sentinel-1C/1D} و ساحل ایران کار کرد؛ برتری
ظاهری بزرگ آن بر \lr{CFAR} نزدیک ساحل به نوار یک‌کیلومتری‌ای برمی‌گردد که \lr{GFW} در آن
گزارشی ندارد، و دورتر از ساحل برتری کوچک اما معنادار می‌ماند.
\item \textbf{شناورهای تاریک:} تلفیق خانه-ساعت با حضور \lr{AIS} نرخ نامزد تاریک را در دوره‌ها و
منطقه‌های گوناگون با روشی یکسان می‌سنجد و هر نامزد را با شواهدش ثبت می‌کند.
\item \textbf{سامانه:} زنجیره‌ای از خواندن تصویر تا کنسول وب فارسی و راست‌به‌چپ، که هر
آشکارسازی را با تصویر راداری‌اش نشان می‌دهد.
\end{itemize}

\section{ساختار گزارش}
فصل~\ref{ch:background} مبانی رادار، \lr{AIS} و پیشینه آشکارسازی شناور را مرور می‌کند.
فصل~\ref{ch:method} داده‌ها، برچسب‌گذاری خودکار، مدل، روش سنجش و تلفیق با \lr{AIS} را شرح
می‌دهد. فصل~\ref{ch:results} نتایج را می‌آورد و فصل~\ref{ch:conclusion} به بحث، محدودیت‌ها و
کارهای آینده می‌پردازد.
